%! Author = lili
%! Date = 7/9/26

%% 1

\newglossaryentry{wann_poisson}{
    name={Wann ist die Poisson-Approximation ist gut?},
    description={Wenn n groß im Vergleich zu k ist ($n >> k$) und p klein ist (
    $n >> \alpha := np$). Das heißt insbesondere, daß die Erfolgswahrscheinlichkeit klein sein muss.},
    type=dinge
}

\newglossaryentry{def_elementarereignis}{
    name={Was ist ein Elementarereignis?},
    description={Ein möglicher Ausgang eines Zufallsexperiments.
    Notation: kleiner griechischer Buchstabe $\omega$ (omega).
    Je nach Kontext auch andere Symbole, z.B. $k$ für natürliche Zahlen oder $x$ für reelle Zahlen.},
    type=dinge
}

\newglossaryentry{def_ereignisraum}{
    name={Was ist ein Ereignisraum $\Omega$?},
description={Die Menge aller Elementarereignisse, also aller möglichen Ausgänge eines Zufallsexperiments.
Muss nichtleer sein ($\Omega \neq \emptyset$). Notation: großer griechischer Buchstabe $\Omega$ (Omega).
Bei mehreren Ereignisräumen: $\Omega, \tilde\Omega, \Omega_1, \Omega_2, \dots$},
type=dinge
}

\newglossaryentry{bsp_ereignisraum_wahl}{
name={Wie wählt man den Ereignisraum beim Würfeln bzw. bei der Zeitmessung?},
description={Würfeln: $\Omega = \{1,\dots,6\}$ (oder Würfelsymbole). Zeitmessung: $\Omega = [0,\infty)$ als sinnvolle Grundwahl, kann bei Zusatzinformationen auf ein kleineres Intervall eingeschränkt werden. Wichtig: Bei Angabe von $\Omega$ stets die Interpretation der Elementarereignisse angeben, z.B. ``Dabei bedeutet $\omega = 1$, dass der Würfel eine Eins zeigt''.},
type=dinge
}

\newglossaryentry{def_ereignis}{
name={Was ist ein Ereignis $A$?},
description={Eine (messbare) Teilmenge des Ereignisraums, $A \subseteq \Omega$. Auch $A=\emptyset$ und $A=\Omega$ sind zulässig. Sprechweise: $\omega \in A$ bedeutet ``$A$ ist eingetroffen''. Formal ist $A$ ein Element der Potenzmenge $\mathcal{P}(\Omega)$, also $A \in \mathcal{P}(\Omega)$.},
type=dinge
}

\newglossaryentry{bem_einelementige_menge}{
name={Ist ein Elementarereignis auch ein Ereignis?},
description={Ja. Ist $\omega \in \Omega$ ein Elementarereignis, so ist die einelementige Menge $\{\omega\} \in \mathcal{P}(\Omega)$ ein Ereignis in $\Omega$.},
type=dinge
}

\newglossaryentry{einschraenkung_ueberabzaehlbar}{
name={Welche Ereignisse betrachtet man, wenn $\Omega \subseteq \mathbb{R}^n$ weder endlich noch abzählbar unendlich ist?},
description={Man beschränkt sich auf Ereignisse, die Intervalle oder bekannte geometrische Figuren sind (Dreieck, Rechteck, Kugel), sowie abzählbare Schnitte oder Vereinigungen solcher Objekte. Dies vermeidet Probleme der Maß- und Integrationstheorie.},
type=dinge
}

\newglossaryentry{def_zaehldichte}{
name={Was ist eine Zähldichte / Wahrscheinlichkeitsfunktion $\mu$?},
description={Sei $\Omega \neq \emptyset$ abzählbar. Eine Abbildung $\mu: \Omega \to [0,1]$ mit $\sum_{\omega \in \Omega} \mu(\omega) = 1$ (Normierung) heißt Zähldichte bzw. Wahrscheinlichkeitsfunktion auf $\Omega$.},
type=dinge
}

\newglossaryentry{def_wahrscheinlichkeitsmass}{
name={Wie berechnet man aus einer Zähldichte $\mu$ die Wahrscheinlichkeit $P(A)$ eines Ereignisses?},
description={$P(A) := \sum_{\omega \in A} \mu(\omega)$. $P$ heißt Wahrscheinlichkeitsmaß und liegt immer zwischen 0 und 1.},
type=dinge
}

\newglossaryentry{konvention_leere_menge}{
name={Welchen Wert hat $P(\emptyset)$?},
description={Konvention: $P(\emptyset) = 0$.},
type=dinge
}

\newglossaryentry{eigenschaften_wmass}{
name={Welche drei Grundeigenschaften hat ein Wahrscheinlichkeitsmaß $P$ auf abzählbarem $\Omega$?},
description={(1) $0 \leq P(A) \leq 1$ für alle $A \in \mathcal{P}(\Omega)$. (2) $P(\emptyset) = 0$. (3) $P(\Omega) = 1$. $P$ ist eine Abbildung $P: \mathcal{P}(\Omega) \to [0,1]$.},
type=dinge
}

\newglossaryentry{beweis_untere_obere_schranke}{
name={Wie zeigt man $0 \leq P(A) \leq 1$?},
description={Untere Schranke: $P(A) = \sum_{\omega \in A} \mu(\omega) \geq 0$, da alle Summanden nichtnegativ sind. Obere Schranke: $P(A) = \sum_{\omega \in A}\mu(\omega) \leq \sum_{\omega \in \Omega}\mu(\omega) = 1$, ebenfalls weil alle Summanden nichtnegativ sind.},
type=dinge
}

\newglossaryentry{def_wraum}{
name={Was ist ein Wahrscheinlichkeitsraum?},
description={Das Paar $(\Omega, \mu)$ heißt Wahrscheinlichkeitsraum. Es besteht eine ein-eindeutige Zuordnung zwischen der Wahrscheinlichkeitsfunktion $\mu$ und dem Wahrscheinlichkeitsmaß $P$: $P(A) = \sum_{\omega \in A}\mu(\omega)$ und umgekehrt $\mu(\omega) = P(\{\omega\})$.},
type=dinge
}

\newglossaryentry{unterschied_mu_P}{
name={Was ist der Unterschied zwischen $\mu$ und $P$ bezüglich ihres Definitionsbereichs?},
description={$\mu$ hat als Argument ein einzelnes Elementarereignis $\omega$ (also $\mu(\omega)$), während $P$ auf der Potenzmenge $\mathcal{P}(\Omega)$ definiert ist und daher das einelementige Ereignis $\{\omega\}$ als Argument hat (also $P(\{\omega\})$).},
type=dinge
}


\newglossaryentry{take_home_modellieren}{
name={Was ist die Take-Home-Message beim Modellieren von $\mu$?},
description={Wann immer möglich, modelliere so, dass $\mu(\omega) = \text{const} \ \forall \omega \in \Omega$ gilt (Gleichverteilung). Dies erfordert Symmetrieargumente.},
type=dinge
}

\newglossaryentry{def_gleichverteilung}{
name={Wann heißt $\mu$ Gleichverteilung auf $\Omega$?},
description={Sei $\Omega \neq \emptyset$ endlich, d.h. $0 < \#\Omega < \infty$. Gilt $\mu(\omega) = \mu(\tilde\omega)$ für alle $\omega,\tilde\omega \in \Omega$ (d.h. $\mu$ ist konstant), so heißen $\mu$ und das zugehörige $P$ die Gleichverteilung auf $\Omega$.},
type=dinge
}

\newglossaryentry{lem_gleichverteilung_rechnen}{
name={Wie berechnet man $\mu(\omega)$ und $P(A)$ bei Gleichverteilung?},
description={(a) $\mu(\omega) = \frac{1}{\#\Omega}$ für alle $\omega \in \Omega$. (b) $P(A) = \frac{\#A}{\#\Omega}$ für alle $A \subseteq \Omega$. Anschaulich: Anzahl der günstigen Ausgänge geteilt durch Anzahl aller möglichen Ausgänge.},
type=dinge
}

\newglossaryentry{beweis_gleichverteilung}{
name={Wie beweist man $\mu(\omega) = \frac{1}{\#\Omega}$ bei Gleichverteilung?},
description={Da $\mu$ konstant ist, existiert ein $c \in [0,1]$ mit $\mu(\omega) = c$ für alle $\omega$. Aus der Normierung folgt $1 = \sum_{\omega \in \Omega}\mu(\omega) = c \cdot \#\Omega$, also $c = \frac{1}{\#\Omega}$.},
type=dinge
}

\newglossaryentry{bem_gleichverteilung_abzaehlprobleme}{
name={Worauf reduziert sich das Berechnen von Wahrscheinlichkeiten bei Gleichverteilung?},
description={Auf reine Abzählprobleme (Kombinatorik), da $P(A) = \frac{\#A}{\#\Omega}$ gilt.},
type=dinge
}

\newglossaryentry{keine_gleichverteilung_unendlich}{
name={Warum kann es keine Gleichverteilung auf einer abzählbar unendlichen Menge geben?},
description={Annahme: Gleichverteilung $\mu$ auf $\Omega = \mathbb{N}$ mit Wert $c$. Aus $1 = \sum_{k=1}^\infty \mu(k) = c \cdot \#\mathbb{N}$ folgt ein Widerspruch: Für $c=0$ ergibt sich $1=0$, für $c>0$ ergibt sich $1=\infty$. Beide Fälle sind unmöglich, also existiert keine Gleichverteilung auf unendlichen Mengen. Deshalb ist in der Definition $\#\Omega < \infty$ vorausgesetzt.},
type=dinge
}

